% Fragmento listo para integrar en la discusión de P4.
% La figura asociada se encuentra en p4-solution.pdf y muestra únicamente
% campos interpolados desde los grados de libertad FEM resueltos.
\subsubsection{P4: solución FEM del benchmark Navier--Stokes DFG 2D-1}

Se resolvió el problema estacionario de Navier--Stokes incompresible en
\(\Omega=([0,2.2]\times[0,0.41])\setminus\overline{B_{0.05}(0.2,0.2)}\), con
densidad \(\rho=1\) y viscosidad cinemática \(\nu=10^{-3}\). La entrada se
impuso como
\[
 \mathbf u_{\mathrm{in}}(y)=\left(\frac{4(0.3)y(0.41-y)}{0.41^2},0\right),
 \qquad U_{\max}=0.3,\quad \overline U=0.2,
\]
las paredes y el cilindro tuvieron no deslizamiento, y en la salida se dejó la
condición natural exacta \(\nu\,\partial_n\mathbf u-p\mathbf n=0\), sin fijar
presión adicional. Así, el número de Reynolds de referencia es
\(\mathrm{Re}=\overline U(0.1)/\nu=20\). El sistema resuelto fue
\[
 -\nu\Delta\mathbf u+(\mathbf u\cdot\nabla)\mathbf u+\nabla p=\mathbf0,
 \qquad \nabla\cdot\mathbf u=0.
\]

La discretización empleó elementos triangulares Taylor--Hood \(P_2/P_1\) para
velocidad y presión. Cada nivel se generó con Delaunay y una frontera
poligonal explícita para el cilindro; el número de segmentos de esa frontera
creció junto con la resolución de fondo. La convección se continuó mediante
Picard y se remató con correcciones de Newton. La tabla resume la corrida
completa que produjo los campos de la figura.

\begin{center}
\small
\begin{tabular}{lrrrrrr}
\hline
nivel & $n_x\times n_y$ & segmentos & elementos & $N_u$ & $N_p$ & Picard+Newton\\
\hline
grueso & $34\times8$ & 32 & 734 & 3168 & 425 & 18+1\\
medio  & $52\times12$ & 48 & 1526 & 6456 & 851 & 18+1\\
fino   & $76\times18$ & 64 & 3100 & 12904 & 1676 & 18+1\\
\hline
\end{tabular}
\end{center}

En el nivel fino, el residuo débil libre fue
\(2.11\times10^{-18}\), el residuo de continuidad débil
\(1.95\times10^{-18}\), el error de traza de entrada
\(4.14\times10^{-17}\) y el máximo de no deslizamiento
\(7.22\times10^{-16}\). El cierre integral de flujo fue
\(2.54\times10^{-15}\) relativo, con flujos de entrada y salida
\(-0.0820\) y \(0.0820\), respectivamente; la traza de salida evaluada con
la condición natural tuvo norma \(7.18\times10^{-7}\). La norma \(L^2\) de
la divergencia evaluada dentro de los elementos bajó de
\(1.40\times10^{-1}\) a \(6.21\times10^{-2}\); la restricción débil de
continuidad se reporta por separado porque es la que impone el espacio mixto.

Los funcionales obtenidos en las tres mallas fueron
\[
\begin{array}{c|rrr}
 & C_D & C_L & \Delta p\\\hline
\text{gruesa} & 5.51699 & 0.00544 & 0.11364\\
\text{media}  & 5.52807 & 0.01015 & 0.12186\\
\text{fina}   & 5.52153 & -0.00651 & 0.11555
\end{array}
\]
con \(\Delta p=p(0.15,0.2)-p(0.25,0.2)\) y
\(\boldsymbol\sigma=\nu\nabla\mathbf u-pI\) integrada con la normal
exterior al cilindro. Como control externo, el benchmark publica
\(C_D=5.57953523384\), \(C_L=0.010618948146\) y
\(\Delta p=0.11752016697\)~\cite{dfgRe20}. La comparación entre niveles da
 cambios relativos medio--fino de \(0.24\%\) en velocidad y \(0.99\%\) en
 presión de media cero, pero \(5.18\%\) en \(\Delta p\); además, el signo de
 \(C_L\) no se estabiliza. Por ello estos resultados documentan una solución
 FEM reproducible y sus balances, pero no se presentan como certificación de
 convergencia de todas las cantidades integrales. Las imágenes de velocidad,
 presión y líneas de corriente se obtuvieron exclusivamente interpolando el
 campo resuelto y enmascarando el mismo polígono sólido de la malla.
